\documentclass[a4paper,11pt]{article}

\usepackage{amsmath, amssymb, amstext, amsfonts, mathrsfs}


% \sffamily %schrift ohne Serifen

\usepackage[T1]{fontenc} 
% schriftencodierung f�r umlaute, trennung
% f\"ur Uni
\usepackage[latin1]{inputenc}
\usepackage{selinput}
% \usepackage[utf8x]{inputenc} 
\usepackage{bibgerm} 
% german bibliography
\usepackage[german]{babel}
%wichtig f�r deutschen Content
\usepackage{ucs}
%erweiterte UTF-8 Unterst�tzung
\usepackage{wrapfig} 
% Paket zur Positionierung einbinden
\usepackage{multirow}
% zusammenfassen von Tabellenzellen
% \usepackage{subscript}
% zum tiefstellen
\usepackage{lscape}
\usepackage{pdflscape}
% zum drehen der Seite
% \usepackage[super]{natbib}
\usepackage[square,sort,comma,numbers]{natbib}
% Erstellung es Literaturverzeichnisses
\usepackage{url}
% Umbruch f�r URL
\usepackage{pst-3dplot}
% f�r tex Grafiken n�tig
\usepackage{pstricks}
\usepackage{listings}
% f�r das einf�gen von Quelltext
\definecolor{codegray}{rgb}{0.92,0.92,0.92}
\lstset{basicstyle=\fontsize{9}{11}\selectfont\ttfamily, breaklines=true, backgroundcolor=\color{codegray}, numbers=left, numberstyle=\tiny, tabsize=4, language=java}
\definecolor{mymauve}{rgb}{0.58,0,0.82}
\definecolor{mygreen}{rgb}{0,0.6,0}
\lstset{
commentstyle=\color{mygreen},
keywordstyle=\color{mymauve},
language=Java,
stringstyle=\color{blue}
}


%Einstellungen f�r Quellcode
\usepackage[a4paper, left=3cm, right=2cm, top=2cm]{geometry}
% Formatierung R�nder
\usepackage[section]{placeins}
% f�r Floatbarriere
\usepackage{color}
\usepackage{colortbl}
%f�r die Verwendung von Farben

\clubpenalty = 10000 
\widowpenalty = 10000
\displaywidowpenalty = 10000
%Verhinderung von Hurenkindern und Schusterjungen
%10000 bedeutet die sie sollen kommplett vermieden werden

\title{Konzept und prototypische Implementierung eines �bergreifenden Dokumenten- und Medienmanagements als SOA-Anwendung}

\author{Sebastian Rieger}

\pagenumbering{arabic}
%Seitenzahlen(arabische Zahlen)

\setlength{\parindent}{0.25cm} 
%Absatzeinzug �ndern in Zoll
\setlength{\parskip}{0.25cm}
%Absatzabstand
\linespread {1.5}
%Zeilenabstand

\usepackage{setspace}
\usepackage{hyperref}
%anklickbare Hyperlinks

%funktioniert nicht bei Fu�noten
\usepackage{graphicx}
\usepackage{graphics}
%f�r einbinden von Grafiken

\usepackage{framed}
%f�r Umrandung der Erkl�rung
\usepackage{acronym}
% f�r abk�rzungen
% \usepackage{PSTricks}
\usepackage{epstopdf}
% f�r eps bilder nutze pdflatex --shell-escape this.tex
\usepackage{amssymb}
% f�r mathematische Symbole

\usepackage{hyperref}
% klickbare links

% \usepackage{pdfpages}
\usepackage{rotating}
\usepackage{svg}
%%%%%%%%%%%%%%%%%%%%%%%%%%%%%%%%%%%%%%%%%%%%%%%%%%%%%%%%%%%%%%%%%%%%%%%%%%%%%%%%%%%%%%%%%%%%%%%%%%%%%
%% Angaben zur Arbeit
%%%%%%%%%%%%%%%%%%%%%%%%%%%%%%%%%%%%%%%%%%%%%%%%%%%%%%%%%%%%%%%%%%%%%%%%%%%%%%%

\newcommand{\Autor}{Sebastian Rieger}
\newcommand{\MatrikelNummer}{7406886}
\newcommand{\Kursbezeichnung}{TINF12B1}

\newcommand{\FirmenName}{Karlsruher Institut f�r Technologie (KIT)}
\newcommand{\FirmenStadt}{Karlsruhe}
\newcommand{\FirmenLogoDeckblatt}{{\includegraphics[width=3cm]{Bilder/kitlogo}}}

% Falls es kein Firmenlogo gibt:
%  \newcommand{\FirmenLogoDeckblatt}{}

\newcommand{\BetreuerFirma}{Dipl. -Inform. Thorsten Schlachter}
\newcommand{\BetreuerDHBW}{Dipl. -Inform., MBA Christian Jordan}
\newcommand{\Titel}{Konzept und prototypische Implementierung eines �bergreifenden Dokumenten- und Medienmanagements als SOA-Anwendung}
\newcommand{\AbgabeDatum}{24.08.2015}

\newcommand{\Dauer}{12 Wochen}

% \newcommand{\Abschluss}{Bachelor of Engineering}
\newcommand{\Abschluss}{Bachelor of Science}

\newcommand{\Studiengang}{Angewandte Informatik}
% \newcommand{\Studiengang}{Angewandte Informatik}
\newcommand{\Was}{Bachelorarbeit}

%%%%%%%%%%%%%%%%%%%%%%%%%%%%%%%%%%%%%%%%%%%%%%%%%%%%%%%%%%%%%%%%%%%%%%%%%%%%%%%%%%%%%%%%%%%%%%%%%%%%% 
%steuervariableISO 23081
\usepackage{ifthen} %Package f�r if/else
\newboolean{bilder} %Deklaration
\setboolean{bilder}{true} %Zuweisung
% \setboolean{bilder}{false} %Zuweisung
%%%%%%%%%%%%%%%%%%%%%%%%%%%%%%%%%%%%%%%%%%%%%%%%%%%%%%%%%%%%%%%%%%%%%%%%%%%%%%%%%%%%%%%%%%%%%%%%%%%%%

\begin{document}

\begin{center}
\vspace*{-2cm}
\FirmenLogoDeckblatt\hfill\includegraphics[width=4cm]{Bilder/dhbw-logo}\\[1cm]
{\Huge \Titel}\\[2cm]
{\Huge\scshape \Was}\\[2cm]
{\large f�r die Pr�fung zum}\\[0.5cm]
{\Large \Abschluss}\\[0.5cm]
{\large des Studienganges \Studiengang}\\[0.5cm]
{\large an der}\\[0.5cm]
{\large Dualen Hochschule Baden-W�rttemberg Karlsruhe}\\[0.5cm]
{\large von}\\[0.5cm]
{\large\bfseries \Autor}\\[1cm]
{\large Abgabedatum \AbgabeDatum}
\vfill
\end{center}
\begin{tabular}{l@{\hspace{1cm}}l}
Bearbeitungszeitraum             & \Dauer                       \\
Matrikelnummer                   & \MatrikelNummer              \\
Kurs                             & \Kursbezeichnung             \\
Ausbildungsfirma                 & \FirmenName                  \\
                                 & \FirmenStadt                 \\
Betreuer der Ausbildungsfirma    & \BetreuerFirma               \\
Gutachter der Studienakademie    & \BetreuerDHBW                \\
\end{tabular}

\newpage
%Seitenumbruch
%%%%%%%%%%%%%%%%%%%%%%%%%%%%%%%%%%%%%%%%%%%%%%%%%%%%%%%%%%%%%%%%%%%%%%%%%%%%%%
%% Descr:       Vorlage für Berichte der DHBW-Karlsruhe, Erklärung
%% Author:      Prof. Dr. Jürgen Vollmer, vollmer@dhbw-karlsruhe.de
%% $Id: erklaerung.tex,v 1.2 2010/07/22 13:30:27 vollmer Exp $
%%%%%%%%%%%%%%%%%%%%%%%%%%%%%%%%%%%%%%%%%%%%%%%%%%%%%%%%%%%%%%%%%%%%%%%%%%%%%%%

% In Bachelorarbeiten muss eine schriftliche Erklärung abgegeben werden. In allen anderen
% Arbeiten entf�llt diese. Hierin best�tigen die Studierenden, dass die Bachelorarbeit
% selbst�ndig verfasst und s�mtliche Quellen und Hilfsmittel angegeben sind. Diese Erkl�rung
% bildet das zweite Blatt der Arbeit. Der Text dieser Erkl�rung muss auf einer separaten Seite
% wie unten angegeben lauten.

\newpage
\thispagestyle{empty}
\begin{framed}
\begin{center}
\Large\bfseries Erkl\"arung
\end{center}

\noindent
Gem\"a\ss{}�16 (3) der "`Studien- und Pr\"ufungsordnung f\"ur den Studienbereich
Technik"' vom 1.11.2007.

\medskip
\noindent
Ich habe die vorliegende Arbeit selbstst\"andig verfasst und
keine anderen als die angegebenen Quellen und Hilfsmittel verwendet.

\vspace{3cm}
\noindent
\underline{\hspace{4cm}}\hfill\underline{\hspace{6cm}}\\
Ort~~~~~Datum\hfill Unterschrift\hspace{4cm}
\end{framed}

%%%%%%%%%%%%%%%%%%%%%%%%%%%%%%%%%%%%%%%%%%%%%%%%%%%%%%%%%%%%%%%%%%%%%%%%%%%%%%%
\endinput
%%%%%%%%%%%%%%%%%%%%%%%%%%%%%%%%%%%%%%%%%%%%%%%%%%%%%%%%%%%%%%%%%%%%%%%%%%%%%%%

\newpage
\begin{spacing}{0.9}

%Einf�gen Inhaltsverzeichnis
\tableofcontents



\end{spacing}

\newpage
\section{Abstract}
\section{Einleitung}
\section{Anforderungen}
\section{Analyse Datenbest\"ande}
\subsection{FADO und Untergruppe}
\subsection{DRS}
\subsection{Bilddatenbank}
\section{Stand der Technik / existierende Konzepte}
\subsection{Fachlich}
\subsection{Technisch}
\section{Erstellung eines Datenkonzepts}
\subsection{Fachlich}
\subsection{Technisch}
\section{Technologievergleich}
\subsection{Alfresco}
\subsection{anderes ECM}
\subsection{anderes ECM}
\section{Implementierung}
\subsection{Schnittstelle}
\subsection{Backend}
\subsection{Frontend}


\newpage
\section{Abk�rzungsverzeichnis}
\begin{acronym}
%   \acro{SDK}{\emph{Software Development Kits}}
%   \acro{OHA}{\emph{Open Handset Alliance}}
%   \acro{IDE}{\emph{Integrierte Entwicklungsumgebung}}
%   \acro{ADT}{\emph{Android Development Tools}}
%   \acro{DVM}{\emph{Dalvik Virtual Machine}}
%   \acro{FME}{\emph{Funkmeldeempf\"anger}}
%   \acro{JIT}{\emph{Just in Time}}
%   \acro{ART}{\emph{Android Runtime}}
%   \acro{AOT}{\emph{Ahead-of-Time-Decodierung}}
%   \acro{API}{\emph{Application Programming Interface}}
  \acro{ECM}{\emph{Enterprise Content Management}}
  \acro{DMS}{\emph{Dokumentenmanagementsystem}}
  \acro{PoC}{\emph{Proof of Concept}}
  \acro{LUBW}{\emph{Landesanstalt f\"ur Umwelt, Messungen und Naturschutz Baden-W\"urttemberg}}
  \acro{GAA}{\emph{Gewerbeaufsicht Baden-W\"urttemberg}}
  \acro{REST}{\emph{Representational State Transfer}}
  \acro{FADO}{\emph{Fachdokumente Online}}
  \acro{DRS}{\emph{Document Retrieval System}}
  \acro{INSPIRE}{\emph{Infrastructure for Spatial Information in the European Community}}
  \acro{API}{\emph{Application Programming Interface}}
  \acro{EU}{\emph{Europ\"aische Union}}
  \acro{}{\emph{}}
%  \acro{KIT}{\emph{Karlsruher Instituts f�r Technologie}}
%  \acro{GDS}{\emph{Generic Data Services}}
% %  \acro{LSDF}{\emph{Large Scale Data Facility}}
%  \acro{OPM}{\emph{Objektorientierten Programmiermodell}}
%  \acro{SMD}{\emph{Strukturelle Metadaten}}
%  \acro{JSON}{\emph{JavaScript Object Notation}}
% %  \acro{HALO}{\emph{High Altitude and Long Range Research Aircraft}}
%  \acro{IAI}{\emph{Institut f�r Angewandte Informatik}}
%  \acro{JAXB}{\emph{Java Architecture for XML Binding}}
%  \acro{UDDE}{\emph{User Data Description Editor}}
%  \acro{AMD}{\emph{Anwendermetadaten}}
%  \acro{CG}{\emph{Class Generator}}
%  \acro{IG}{\emph{Interface Generator}}
\end{acronym}
\newpage
\listoffigures
\newpage
\listoftables
% Abk�rzungsverzeichnis
\newpage
\bibliographystyle{alpha}
% verzeichnis im DIN format
\bibliography{Quellen}
\end{document}
